% =============================================================================
\chapter{A Survey of Agents for Computer Use Hands Two Gaps to This Thesis}
\label{chap:acu}
\margintoc
% =============================================================================
%
% Thesis owned throughout. The survey is summarised rather than reproduced, so
% this chapter opens no importedpaper block and the paper's appendix is not in
% the thesis. imports/acu.toml carries the decision and the reason. Wiring the
% paper back in means setting status = "wired" and restoring the includes.
%
% The six gaps are written from imported_papers/ACU_JAIR/sections/09_conclusion.tex
% and nowhere else. The version of this chapter that reproduced the paper
% embellished them past what the survey says. 87 is the number of agents
% classified, not of papers read: the paper count is larger. The survey says
% "we classify \numagents as ACU agents and \numdatasets as computer use
% datasets".

% ─────────────────────────────────────────────────────────────────────────────
% PUBLICATION HEADER
% ─────────────────────────────────────────────────────────────────────────────

\begin{publicationbasis}{Published}
Sager, P.\,J., Meyer, B., Yan, P., von Wartburg-Kottler, R., Etaiwi, L.,
Enayati, A., Nobel, G., Abdulkadir, A., Grewe, B.\,F., \& Stadelmann, T.
(2026). A comprehensive survey of agents for computer use: Foundations,
challenges, and future directions. \textit{Journal of Artificial Intelligence
Research}, \textit{85}, article 34. \cite{sager_survey_2026}
\end{publicationbasis}

% ─────────────────────────────────────────────────────────────────────────────
\contributionstatement
\label{sec:acu_contribution}
% ─────────────────────────────────────────────────────────────────────────────

% TODO(author): the four paragraphs below were written in an earlier session
% and state the author's own attribution, which is the author's to make. They
% are left as written. One reviewer found the register advocative in places,
% in particular "initiated the survey independently" and the two separate
% statements that attribution is unclear.

The thesis author holds shared first authorship on this work with Benjamin
Meyer. The vendored manuscript marks the thesis author as corresponding author
and carries its equal contribution note commented out, so the shared first
authorship has to be checked against the published article before
submission. The collaboration was iterative and the contributions of both first
authors are interleaved throughout the paper in a way that makes clean
attribution of individual sections difficult. The following account reflects
the thesis author's best reconstruction of how the work developed.

The project originated with the thesis author, who identified the absence of a
theoretically grounded taxonomy that did not depend on the technology as a gap
in the existing ACU literature, and initiated the survey independently. The
thesis author proposed the initial structure of the taxonomy in a first draft,
together with the selection of the three analytical perspectives: domain,
interaction, and agent. Benjamin Meyer subsequently refined the taxonomy
structure substantially, improving its internal consistency and the precision
of its category boundaries, after which the thesis author revised the framework
again. The final taxonomy came out of several such passes and cannot be
attributed to either first author alone.

The paper collection and classification process was carried out collectively by
all authors: each author independently read and classified a subset of the
surveyed papers into a shared structured spreadsheet, using the evolving
taxonomy as a classification schema. The survey structure, including the
identification and articulation of the six research gaps, was derived
inductively from the classification data through discussions involving both
first authors and with input from the senior authors Ahmed Abdulkadir,
Benjamin Grewe, and Thilo Stadelmann.

The reading of the survey in this chapter belongs to the thesis and not to
the paper. It sorts the six gaps by whether they are about what training produced,
and it reads the first of them in a way the survey raises only in passing,
which the sentence that does so says. This paper is summarised, as the four in
\cref{chap:further} are, and it keeps a chapter of its own because the
authorship here is shared first where the contribution to those four was
partial, and because the thesis takes its question from this work.

% ─────────────────────────────────────────────────────────────────────────────
% \texorpdfstring, because hyperref strips math from a PDF bookmark and warns
% about it. 3.3 wants the number in math mode in the printed heading.
\section{The survey classifies \texorpdfstring{$87$}{87} agents and excludes the best known commercial systems}
\label{sec:acu_scope}
% ─────────────────────────────────────────────────────────────────────────────

Agents that operate a computer from an instruction in natural language reach
success rates almost six times lower than people. The survey cites three
causes: the interface changes while the agent is working, the task requires
switching between applications, and the interface is dense, as a spreadsheet
is. The figure comes from one benchmark, which the survey cites rather than
runs~\cite{sager_survey_2026}. This thesis starts at that distance between
what these systems are asked to do and what they manage.

The survey classifies $87$ agents and $33$ datasets from publications between
$2018$ and $2024$. The number of papers read is larger, because agents and
datasets are separated in the analysis and a paper may supply either.

The survey draws its boundary in five places, each for a reason. Agents that only give
advice are out, because the survey is about systems that act. Agents that play
video games, control server facilities, write code or test software are out,
because each is a field with its own literature. Only instructions given as
text are covered, and only agents built on deep learning, which leaves out the
rule based automation that does much of this work in practice. Coverage is
also restricted to fully disclosed systems, which leaves out some of the
recent commercial products, the ones a reader in $2026$ thinks of first.

The survey is not a systematic review, and says so. The field lacks a shared
vocabulary, which makes a classic systematic review infeasible, so the
collection was iterative, combining expert knowledge with backward and forward
snowball sampling until no further relevant work appeared. Two researchers
screened each candidate independently and resolved disagreements in
discussion. About a third of the cited works are preprints, and the survey
says these may lack validation and should be read with that in mind.

% ─────────────────────────────────────────────────────────────────────────────
\section{A taxonomy built on agent theory brings two literatures into one scheme}
\label{sec:acu_taxonomy}
% ─────────────────────────────────────────────────────────────────────────────

The survey organises the literature along three perspectives. The domain
perspective characterises the computing environment the agent works in. The
interaction perspective describes what the agent observes and what actions it
can take. The agent perspective concerns the internal structure of the agent
and how it learns.

Building the scheme on agent theory, and on no particular technology, brings
work together that had never been compared. Earlier surveys each cover one
technology or one platform: some review only agents built on foundation models
and leave out reinforcement learning as a design principle, others review
reinforcement learning agents in general without treating computer use, and
one covers mobile agents alone. One earlier survey has a similar scope, and
the paper's claim against it rests on depth. A scheme that applies both to an
agent trained by reinforcement learning and to an agent assembled by prompting
lets the whole area be read as one field.

The paper states four limits of the taxonomy. It supports no comparison of
capability across agents, because many of the surveyed systems cover only part
of a benchmark and report different metrics, which makes a common ranking
infeasible. It carries no analysis of how individual design choices affect
performance. It covers instructions given as text. The fourth limit concerns
an observation and action loop that changes over time, which the scheme would
accommodate while the survey applies it only to interaction patterns that stay
fixed, a choice of scope the scheme itself does not force.

% ─────────────────────────────────────────────────────────────────────────────
\section{The field is moving to foundation agents and has not settled how they learn}
\label{sec:acu_field}
% ─────────────────────────────────────────────────────────────────────────────

Specialised agents built for one narrow task are giving way to agents built on
foundation models. The survey reads that shift as a consequence of what open
ended reasoning over a general computer demands of a model. Specialised
designs remain viable where the task is narrow and efficiency matters.

How those agents learn is not settled, and the survey says in as many words
that the field has yet to converge on a unified framework. Counting agents by
strategy, $31$ clone human demonstrations, $12$ use reinforcement learning,
$12$ keep a long term memory and $30$ learn nothing from the environment at
all. Those figures are the survey's own and they total $85$ against the $87$
agents, so they count strategies and leave the agents unpartitioned. Among the
agents that do learn from experience, behaviour cloning is the most common
choice.

The modality the agents observe through is moving in one direction without
having arrived. Counting agents by modality, $35$ observe through text alone, $31$ through
images, $15$ through both and $4$ indirectly, which again totals $85$ rather
than $87$. Text is still the most common single modality. The movement is in the yearly counts: agents observing through
images go from none in $2022$ to $12$ in $2023$ and $19$ in $2024$, while the
text counts over the same years are $3$, $10$ and $8$.

% ─────────────────────────────────────────────────────────────────────────────
\section{Six gaps recur across the classified agents}
\label{sec:acu_gaps}
% ─────────────────────────────────────────────────────────────────────────────

Applying the taxonomy to the $87$ agents yields six limitations that recur across them.

\begin{enumerate}
    \item \textbf{Insufficient generalisation.} Agents rely on input
    modalities whose structure is inconsistent across applications, which
    hinders generalisation. The survey's own diagnosis is the modality:
    structured text such as a page's markup or a view hierarchy carries
    semantic information that is optional, and is often absent, incomplete,
    inconsistent or ambiguous in real settings.

    \item \textbf{Inefficient learning.} The survey finds no effective and
    practical way for a foundation agent to learn from its environment.
    Reinforcement learning and behaviour cloning both give a strong learning
    signal, and both demand either a simulator of high fidelity or a curated
    and labelled set of demonstrations. A long term memory is practical and
    generalises poorly.

    \item \textbf{Limited planning.} Executing a complex task over many steps
    requires planning, and current agents plan poorly. The survey locates the
    difficulty partly in the model, whose planning over long horizons is weak,
    and partly in the environment, whose dynamics are usually unknown.

    \item \textbf{Low task complexity in benchmarks.} Benchmark designers have
    pursued realistic perception ahead of realistic tasks. The missing
    ingredient is structural and causal dependency between the actions in a
    sequence, which is where the difficulty of a real task lies.

    \item \textbf{Non standardised evaluation.} Metrics differ between papers,
    and success rates in particular are computed in ways that are not
    comparable. Direct comparison between systems is therefore unreliable.

    \item \textbf{A disconnect between research and practical conditions.}
    Agents are built for environments assumed to be deterministic, static and
    episodic, while real environments change under the agent and change again
    when the software is updated. Privacy is a second difficulty, because a
    user cannot control what an agent operating autonomously on their machine
    observes and sends on, and safety a third, because the field studies full
    autonomy while an agent that hands control back for a critical decision
    would be safer.
\end{enumerate}

The first five are discussed throughout the survey's own body. The sixth is
different in kind. The paper presents it as a set of real world discrepancies
that the current literature does not address, which makes it the survey's own
addition to the list.

Against the six the paper sets six directions, which are as much of its
contribution as the gaps are. It recommends observation spaces built on images
for consistent perception, learning strategies that cost less and so scale and
adapt better, policy architectures that support reasoning and planning over
long horizons, benchmarks that combine realistic perception with task
complexity, standardised metrics and success rates in particular, and research
grounded in the conditions under which these agents would actually run.

% ─────────────────────────────────────────────────────────────────────────────
\section{Two of the six gaps are about what training produced}
\label{sec:acu_takes}
% ─────────────────────────────────────────────────────────────────────────────

Those two gaps are where the thesis begins. Three are about evaluation and about the conditions the field assumes when it
builds. The remaining one, planning over many steps, does not sort, for a
reason the fourth paragraph of this section sets out. None of the other four
is this thesis's subject.

Taking the first that way is a reading, and the survey supports it only in
passing. Its own diagnosis of insufficient generalisation is the input
modality, and the remedy it proposes is to change what the agent observes. The
sentence this thesis leans on is the survey's aside that an agent trained on
such representations may develop heuristics specific to the situations it saw,
which do not carry beyond the training data. The survey's own diagnosis
points at the input format, and a heuristic learned in training is a property
of what training produced. The survey offers the point as a hypothesis.

The second gap needs a distinction. The survey counts the cost of the whole
training pipeline, simulators and labelled demonstrations included, and its
recommendation asks for learning strategies that cost less and adapt better,
so it already reaches partway into what follows. The component this thesis
takes is the narrower one, how much data a network needs before it generalises
and how well what it learnt carries to a related task.

Planning is the gap that does not sort. The survey puts the difficulty in a model whose planning over
long horizons is weak and in an environment whose dynamics are unknown, and
its recommendation asks for better policy architectures while one of its two
remedies is a hybrid with symbolic planning. It therefore sits partly inside
what training produced and partly outside it. This thesis measures neither
part.

The survey cannot settle what causes the two gaps it hands over. It classifies $87$ agents by learning strategy and compares
across them, so it does vary the procedure. It never varies the procedure with
the task, the data, the architecture and the width held fixed, which is the
only comparison in which a difference between two trained networks can be
attributed to how they were trained, and no survey could. A success rate
cannot separate two networks that score alike either. The rest of this thesis
therefore measures inside the network, on tasks small enough that the solution
can be read neuron by neuron.

The survey does not predict the result of \cref{chap:grokking}, and no part of
the argument rests on it. The question came from there.
\Cref{chap:further} takes up a paper that sharpens the first of those two
gaps, and another that shows a score missing what it stands for, before
\cref{part:procedure} turns to what decides them.
